\documentclass[11pt]{article}
\usepackage{amsmath,amssymb}
\usepackage[margin=1in]{geometry}
\title{Working equations for the RHF DLPNO-CCSD(T) nuclear Hessian}
\author{Quantum Chemistry Grand Battleplan}
\date{7 November 2026}
\begin{document}
\maketitle

\section{Target and differentiability}
Use the ordinary iterative (T) gradient-model energy with the conventions of
the companion gradient equation record. Indices $i,j,k$ label active occupied
LMOs, $\mu,\nu$ AOs; $A,B$ below index Cartesian nuclear coordinates in bohr.
For a branch $\sigma$ with fixed pair and triplet membership, domain and pivot
lists, and retained PNO/TNO ranks, define
\begin{equation}
 g_A=\partial_A E_\sigma,\qquad H_{AB}=\partial_B g_A,\qquad
 (Hv)_A=\left.\frac{d}{ds}g_A(\mathbf X+s v)\right|_{s=0}.
 \label{eq:target}
\end{equation}
The orbital coefficients, subspace projectors and amplitudes relax. Assume
a twice differentiable solution with separated retained/discarded spaces
and a locally determined response after removing redundant rotations.

\section{Constrained curvature}
Let $q$ denote independent real electronic coordinates after fixing gauges,
$R$ their defining residuals, and $\lambda$ their multipliers. Contractions
in vector notation include the ordered-pair weights of the gradient record.
The scalar Lagrangian and stationary equations are
\begin{equation}
 \mathcal L=E(q,\mathbf X)+\lambda^{\mathsf T}R(q,\mathbf X),\qquad
 R=0,\qquad E_q+R_q^{\mathsf T}\lambda=0 .
\end{equation}
Use column gradients; subscripts on $\mathcal L$ below are partial derivatives
at the reference solution. For dimensionless directions $u,v$ write
$\partial_v=\sum_A v_A\partial_A$. First-order tangents satisfy
\begin{equation}
 R_q q^v=-R_v .
 \label{eq:tangent}
\end{equation}
Since $R$ vanishes along the solution, its contraction with a multiplier held
at the reference value adds zero along the entire local branch. Differentiating
that scalar twice leaves
\begin{align}
 u^{\mathsf T}Hv
 &=\mathcal L_{uv}+\mathcal L_{uq}q^v+(q^u)^{\mathsf T}\mathcal L_{qv}
   +(q^u)^{\mathsf T}\mathcal L_{qq}q^v . \label{eq:curvature}
\end{align}
The coefficient of $q^{uv}$ is $\mathcal L_q=0$.
Alternatively, to obtain the entire column for one tangent $q^v$, solve the
differentiated left stationarity equation and differentiate the analytic gradient:
\begin{align}
 R_q^{\mathsf T}\lambda^v&=-\mathcal L_{qv}-\mathcal L_{qq}q^v,\\
 (Hv)_A&=\mathcal L_{Av}+\mathcal L_{Aq}q^v+R_A^{\mathsf T}\lambda^v .
 \label{eq:column}
\end{align}
Equation~\eqref{eq:column} avoids a separate right solve for each output row.
Multiplying its last term by $u_A$ and summing, then using
Eq.~\eqref{eq:tangent}, recovers Eq.~\eqref{eq:curvature}.
These identities specify the analytic target; local residual contractions
will be expanded as each block is derived and tested.

\section{Canonical verification at one geometry}
Put O at the origin, H1 on $+z$, and H2 in the $yz$ plane with positive $y$.
Use $r=0.958$ angstrom and $\theta=104.4776$ degrees:
\begin{equation}
 \mathbf X_{\rm H1}=r(0,0,1),\qquad
 \mathbf X_{\rm H2}=r(0,\sin\theta,\cos\theta),\quad
 v=e_{{\rm H1},z},\quad u=e_{{\rm O},y}.
\end{equation}
Convert $r$ to bohr before displacement. Hold the laboratory frame fixed.
For an analytic canonical gradient $g$, define
\begin{align}
 C(h;v)&=\frac{g(\mathbf X+h v)-g(\mathbf X-h v)}{2h}
       =Hv+O(h^2),\\
 T(h;v)&=\frac{4C(h/2;v)-C(h;v)}{3}=Hv+O(h^4).
\end{align}
The orders assume sufficient smoothness. Compare $T(0.004;v)$ with
$T(0.008;v)$ in bohr; the former is the reported reference. Check
\begin{equation}
 u^{\mathsf T}Hv=v^{\mathsf T}Hu,\qquad
 \sum_{\text{atoms }K}(Hv)_{K\alpha}=0\quad(\alpha=x,y,z).
\end{equation}
No mass weighting or projection is applied. A directional curvature at
this test geometry has units of hartree/bohr$^2$.

\section{Reusing the restricted SCF nuclear response}
Psi4's restricted Hessian routine forms the nuclear CPHF right-hand sides,
solves the coupled equations, and writes $U_{pi}^{A}$ for all occupied
columns (including frozen columns). Here $C$ contains all reference MOs and
$C_o$ contains all occupied columns. Contract the Cartesian records with
$v_A$ and use the native occupied gauge
\begin{equation}
 U_{ij}^{v}=-\tfrac12(C_o^{\mathsf T}S^v C_o)_{ij},\qquad
 C_o^v=CU^v,\qquad D^v=C_o^v C_o^{\mathsf T}+C_o(C_o^v)^{\mathsf T}.
\end{equation}
The native virtual--occupied record already includes the sign conversion
from Psi4's CPHF unknown to the coefficient derivative. With $D=C_o C_o^{\mathsf T}$,
the projected relaxed Fock derivative (at fixed external MO columns) is
\begin{equation}
 f^v=C^{\mathsf T}\left(F_{\rm explicit}^{v}+2J[D^v]-K[D^v]\right)C_o.
\end{equation}
The native records are Fpi for the explicit derivative and Qpi for the
induced term. For active occupied $i$ and frozen $c$, rotate these blocks by
\begin{align}
 k_{ic}&=-\frac{f^v_{ic}+\epsilon_i U^v_{ic}+\epsilon_c U^v_{ci}}
                   {\epsilon_i-\epsilon_c},&
 U^v_{ic}&\gets U^v_{ic}+k_{ic},&
 U^v_{ci}&\gets U^v_{ci}-k_{ic}.
\end{align}
This maintains zero mixed canonical Fock element to first order. It preserves
$D^v$ and changes the separate core and active coefficient derivatives.
With $C_v$ the virtual columns, the verification conditions are
\begin{align}
 (C_o^v)^{\mathsf T}SC_o+C_o^{\mathsf T}SC_o^v+C_o^{\mathsf T}S^v C_o&=0,\\
 C_v^{\mathsf T}\big[F^v C_o+FC_o^v-(S^v C_o+SC_o^v)\epsilon_o\big]&=0 .
\end{align}

\section{Nuclear tangent of Boys-localized active occupied orbitals}
Let $C_a$ contain the active occupied columns and $L$ the actual forward
Boys LMOs. Let $U=C_a^{\mathsf T}SL$. Along the nuclear direction $v$
define the transported coefficients $T=C_a^v U$. In the final LMO frame,
the unknown additional rotation $K=-K^{\mathsf T}$ gives
\begin{equation}
 L^v=T+LK,\qquad K_{ji}=k_{ij}=-K_{ij}\quad (i<j).
\end{equation}
For each of the three AO electronic dipole operators $d_m$, put
$A_m=L^{\mathsf T}d_m L$. At fixed localization rotation its derivative is
\begin{equation}
 B_m=T^{\mathsf T}d_m L+L^{\mathsf T}d_m T+L^{\mathsf T}d_m^v L,\qquad
 (A_m)^v=B_m+A_mK-KA_m .
\end{equation}
Write $\Delta_{ij}A=A_{ii}-A_{jj}$. Differentiating the objective
$\sum_{mi}(A_m)_{ii}^2$ with respect to $k_{ij}$ gives the stationary
equations and their driving term
\begin{align}
 g_{ij}&=4\sum_m\Delta_{ij}A_m (A_m)_{ij}=0,\\
 b_{ij}^v&=4\sum_m [\Delta_{ij}B_m (A_m)_{ij}
                      +\Delta_{ij}A_m (B_m)_{ij}],\\
 H_{\rm Boys} k&=-b^v .
\end{align}
Here $H_{\rm Boys}$ is the same exponential-coordinate rotation Hessian used
by the gradient adjoint. Require a stationary, locally determined maximum.
For each column $l_i$ verify
\begin{equation}
 (l_i l_i^{\mathsf T})^v=l_i^v l_i^{\mathsf T}+l_i(l_i^v)^{\mathsf T}.
\end{equation}
The rank-one AO density factors are sign-invariant. The coefficient-space
projector acting in the AO metric would have an additional right factor
$S$, which is itself differentiated. Independently converged displaced
solutions are matched by maximizing absolute cross-geometry overlaps.

\section{Raw PAO projector and normalized PAO coefficients}
Let $O=[C_{\rm c},L]$ contain frozen and localized active orbitals. Their
metric is $O^{\mathsf T}SO=I$. With PAO columns indexed by their parent AOs,
\begin{align}
 B &= I-OO^{\mathsf T}S,\\
 B^v &=-O^v O^{\mathsf T}S-O(O^v)^{\mathsf T}S-OO^{\mathsf T}S^v .
\end{align}
For $L^v=T+LK$, the internal rotation contributes exactly
\begin{equation}
 -L(K+K^{\mathsf T})L^{\mathsf T}S=0 .
\end{equation}
Compute $B^v_{\rm loc}$ using the Boys and frozen derivatives, and
$B^v_{\rm can}$ using all canonical occupied coefficients with their RHF
derivatives. Define the diagnostic by the largest elementwise difference,
\begin{equation}
 e_{\rm rot}=\max_{\mu\nu}|(B^v_{\rm loc}-B^v_{\rm can})_{\mu\nu}|.
\end{equation}
The per-column squared norm $m=b^{\mathsf T}Sb$ and normalized coefficient
$n=b/\sqrt m$ have derivatives
\begin{align}
 m^v &=2(b^v)^{\mathsf T}Sb+b^{\mathsf T}S^v b,\\
 n^v &=m^{-1/2}b^v-\tfrac12 b\,m^{-3/2}m^v .
\end{align}
Use the same fixed zero-column branch as the energy for $m\le10^{-20}$.
Differentiated idempotency $B^vB+BB^v=B^v$ and column normalization
$2(n^v)^{\mathsf T}Sn+n^{\mathsf T}S^vn=0$ give additional checks.

\section{First retained PNO space of one pair}
Hold the AO, PAO and fitting domains, parent pivots, AO shell screening
mask and forward retained rank of pair $ij$ fixed. On each fitting center
let $l,r$ denote the left and right AO index lists. For LMO $i$ and the
normalized pair PAO columns $N$, the fitted LMO coefficients and transformed
three-center integrals obey (one auxiliary index $Q$ is implicit)
\begin{align}
 S_{ll}c &= S_{l,:}L_i, &
 S_{ll}c^v&=S_{l,:}^vL_i+S_{l,:}L_i^v-S_{ll}^v c,\\
 q_i&=c^{\mathsf T} I_{lr}N_r,&
 q_i^v&=(c^v)^{\mathsf T}I_{lr}N_r+c^{\mathsf T}I_{lr}^vN_r+c^{\mathsf T}I_{lr}N_r^v .
\end{align}
Stack the rows in pair auxiliary order. With the local RIFIT metric $M$,
\begin{align}
 M b_i=q_i,\quad M b_i^v=q_i^v-M^v b_i,\qquad
 K=b_i^{\mathsf T}q_j,\quad K^v=(b_i^v)^{\mathsf T}q_j+b_i^{\mathsf T}q_j^v .
\end{align}
Use the forward semicanonical PAO transform $X$, padded outside its fixed
pivots. The present tangent requires its nonzero rows to span the full pivot
space. Put $Y=NX$, $W=N^vX$. A metric-preserving frame is
\begin{align}
 G&=-\tfrac12(W^{\mathsf T}SY+Y^{\mathsf T}SW+Y^{\mathsf T}S^vY),&
 Y^v&=W+YG .
\end{align}
Thus $f=Y^{\mathsf T}FY$, $k=X^{\mathsf T}KX$, with derivatives
\begin{align}
 f^v&=(Y^v)^{\mathsf T}FY+Y^{\mathsf T}F^vY+Y^{\mathsf T}FY^v,\\
 k^v&=X^{\mathsf T}K^v X+G^{\mathsf T}k+kG .
\end{align}
At the origin $f_{ab}=\epsilon_a\delta_{ab}$. Write $s=F_{ii}+F_{jj}$
in the active LMO basis, so that $s^v$ includes the LMO derivatives and
the relaxed AO Fock derivative. Differentiating $s t-ft-t f=k$ yields
\begin{equation}
 t_{ab}=\frac{k_{ab}}{s-\epsilon_a-\epsilon_b},\qquad
 t^v_{ab}=\frac{[k^v+f^v t+t f^v-s^v t]_{ab}}{s-\epsilon_a-\epsilon_b}.
\end{equation}
For $\tilde t=2t-t^{\mathsf T}$ and $\eta=(1+\delta_{ij})^{-1}$,
\begin{align}
 D&=\eta(\tilde t t^{\mathsf T}+\tilde t^{\mathsf T}t),\\
 D^v&=\eta(\tilde t^v t^{\mathsf T}+\tilde t (t^v)^{\mathsf T}
              +(\tilde t^v)^{\mathsf T}t+\tilde t^{\mathsf T}t^v).
\end{align}
Let $U_r,U_d$ be the retained and discarded eigenvectors of $D$ with
occupations $n_p,n_a$. In the transported frame,
\begin{align}
 Z_{ap}&=(U_d^{\mathsf T}D^vU_r)_{ap}/(n_p-n_a),\\
 P&=U_r U_r^{\mathsf T},\qquad P^v=U_d Z U_r^{\mathsf T}+U_r Z^{\mathsf T}U_d^{\mathsf T},\\
 A&=YPY^{\mathsf T},\quad
 A^v=Y^v P Y^{\mathsf T}+YP^vY^{\mathsf T}+YP(Y^v)^{\mathsf T}.
\end{align}
$A$ is the AO contravariant subspace density; $AS$ acts on AO coefficients.
Differentiating $ASA=A$ supplies a metric diagnostic.

\section{Coupled first-space LMP2 amplitude response}
For each retained pair $p=(ij)$ let $Y_p$ be its actual first-stage
semicanonical AO frame, $Y_p^{\mathsf T} S Y_p=I$, and $P_p=Y_pY_p^{\mathsf T}$.
From the previous section's AO density derivative $P_p^v$, choose
\begin{equation}
 Y_p^v=P_p^v S Y_p+\tfrac12Y_p(Y_p^{\mathsf T}S^vY_p).
\end{equation}
This fixes the internal gauge, $Y_p^{\mathsf T}S Y_p^v
=-Y_p^{\mathsf T}S^vY_p/2$. Define the following blocks in these frames:
\begin{align}
 f_p&=Y_p^{\mathsf T}F Y_p,& s_p&=F_{ii}+F_{jj},&
 S_{pq}&=Y_p^{\mathsf T}S Y_q,\\
 f_p^v&=(Y_p^v)^{\mathsf T}FY_p+Y_p^{\mathsf T}F^vY_p+Y_p^{\mathsf T}FY_p^v,\\
 S_{pq}^v&=(Y_p^v)^{\mathsf T}S Y_q+Y_p^{\mathsf T}S^v Y_q+Y_p^{\mathsf T}S Y_q^v .
\end{align}
Here $F_{ij}$ is the LMO Fock matrix and $K_p$ the exchange integral
matrix in frame $Y_p$. For each actual directed coupling $p\gets q$ in
the screened forward row, $c_{pq}$ denotes its occupied Fock element.
Use the physical orientation of $q$, transposing the stored unique block
if the row calls for $T_{ji}=T_{ij}^{\mathsf T}$. The right residual is
\begin{equation}
 R_p=K_p+f_pT_p+T_pf_p-s_pT_p-
        \sum_{p\gets q}c_{pq}S_{pq}T_qS_{pq}^{\mathsf T}.
\end{equation}
Let $A$ be its derivative with respect to all independent amplitude blocks.
The moving-basis right-hand side at fixed amplitudes is
\begin{align}
 b_p^v={}&K_p^v+f_p^vT_p+T_pf_p^v-s_p^vT_p \notag\\
 &-\sum_{p\gets q}\big[
 c_{pq}^v S_{pq}T_q S_{pq}^{\mathsf T}
 +c_{pq}(S_{pq}^vT_q S_{pq}^{\mathsf T}
 +S_{pq}T_q(S_{pq}^v)^{\mathsf T})\big],\\
 A T^v={}&-b^v .
\end{align}
The diagonal-pair blocks are symmetric. Compare the gauge-independent AO
tensor derivative in an aligned occupied-pair convention:
\begin{equation}
 (Y_pT_pY_p^{\mathsf T})^v=Y_p^v T_p Y_p^{\mathsf T}
 +Y_p T_p^v Y_p^{\mathsf T}+Y_p T_p(Y_p^v)^{\mathsf T}.
\end{equation}

\section{Final CCSD PNO subspace after coupled LMP2}
For $p=(ij)$ use the first-stage frame $Y$ and its derivative $Y^v$ from
the previous section. Here $T$ is the converged coupled LMP2 amplitude
matrix in $Y$ and $T^v$ its derivative in this same gauge. Define
$\tilde T=2T-T^{\mathsf T}$ and $\eta=(1+\delta_{ij})^{-1}$.
The final density and its derivative in the moving parent frame are
\begin{align}
 D &=\eta(\tilde T T^{\mathsf T}+\tilde T^{\mathsf T}T),\\
 D^v&=\eta\big[\tilde T^v T^{\mathsf T}+\tilde T(T^v)^{\mathsf T}
        +(\tilde T^v)^{\mathsf T}T+\tilde T^{\mathsf T}T^v\big].
\end{align}
For the fixed forward rank $r$, let $P=U_rU_r^{\mathsf T}$ project onto the
$r$ leading eigenvectors of $D$. For retained $b$ and discarded $a$,
\begin{align}
 Z_{ab}&=\frac{(U_d^{\mathsf T}D^vU_r)_{ab}}{n_b-n_a},&
 P^v &= U_d Z U_r^{\mathsf T}+U_r Z^{\mathsf T} U_d^{\mathsf T},\\
 A &=Y P Y^{\mathsf T},&
 A^v &= Y^v P Y^{\mathsf T}+Y P^v Y^{\mathsf T}+YP(Y^v)^{\mathsf T}.
\end{align}
Even when $r$ equals the parent rank, $A^v$ includes parent motion.
The final semicanonical rotation within the retained space drops out of $A$.
The AO coefficient-space projector is $AS$ and its derivative is $A^v S+A S^v$.
As a sign and gauge test, changing $Y^v$ by $YK$ and $T^v$ by
$TK-KT$ for $K=-K^{\mathsf T}$ leaves $A^v$ unchanged.

\section{Starting the CCSD residual source}
For a final CCSD frame $Y_p$, insert the final $A^v$ into the parallel
transport formula used for the first PNO frame. Use
\begin{equation}
 S_{pq}^v=(Y_p^v)^{\mathsf T}SY_q+Y_p^{\mathsf T}S^vY_q+
 Y_p^{\mathsf T}S Y_q^v .
\end{equation}
All residual sources in this section hold the CCSD amplitudes fixed in the
transported frames. The eventual linear equation is $R_{,t}t^v=-b^v$;
the following contributions to $b^v$ are implemented.

\subsection{Projected Fock blocks}
Write $o=L^{\mathsf T}FL$, $h_p=L^{\mathsf T}FY_p$,
$f_p=Y_p^{\mathsf T}FY_p$. For any such block $M=C_l^{\mathsf T}FC_r$,
\begin{equation}
 M^v=(C_l^v)^{\mathsf T}FC_r+C_l^{\mathsf T}F^v C_r+
 C_l^{\mathsf T}FC_r^v .
\end{equation}
Let $n_p$ contain the forward projected singles $S_{p,mm}t_m$ as rows for
$m$ in the pair domain, and $u_p=2t_p-t_p^{\mathsf T}$.
For the partial derivative through bare Fock blocks (hold all explicit
overlap occurrences fixed) define
\begin{align}
 V_p&=f_p^v-n_p^{\mathsf T}h_{p,\mathrm{dom}}^v,\\
 w_{kj}&=o^v_{kj}+(h^v_{jj})_{k:}\,t_j
 \quad\text{on a retained }kj .
\end{align}
Outside the retained list use $w_{kj}=o^v_{kj}$ as in the forward occupied
Fock intermediate. With $p=ij$ and $d=ii$, the sources are
\begin{align}
 b_{1,i}^{F}&=(h_d^v)_{i:}^{\mathsf T}+V_d t_i-
 n_d^{\mathsf T}o^v_{\mathrm{dom},i}
 +\sum_{k} S_{ik,d}^{\mathsf T}u_{ik}(h_{ik}^v)_{k:}^{\mathsf T},\\
 G_{ij}&=-\sum_{k}w_{kj} S_{p,ik}t_{ik}S_{p,ik}^{\mathsf T},\\
 b_{2,ij}^{F}&=t_{ij} V_p^{\mathsf T}+V_p t_{ij}
                  +G_{ij}+G_{ji}^{\mathsf T}.
\end{align}
The sums use the same occupied neighbors as the right residual. The independent
Lambda check compares $2\sum_i (h_{ii}^v)_{i:}t_i+\lambda_1:b_1^F+
\tfrac12\sum_{ij}\lambda_{2,ij}:b_{2,ij}^F$ to the native Fock covectors
contracted with the same perturbation.

\subsection{Direct fitted doubles product and virtual ladder}
On the pair auxiliary domain set $W=M^{-1/2}$ with the same filtered
eigenspectrum as the gradient. For a raw block $q$, $B=Wq$ and
$B^v=W^v q+W q^v$. With $M=U\operatorname{diag}(m)U^{\mathsf T}$,
\begin{equation}
 W^v=U\big[f^{[1]}(m_a,m_b)\circ(U^{\mathsf T}M^v U)\big]U^{\mathsf T},
\end{equation}
where $f^{[1]}$ is the fixed-branch divided difference of the forward
filtered inverse square root (the same self-adjoint linear map used in the
CCSD fitting adjoint). For each auxiliary row, let
$n^v_{m:}=(S^v_{p,mm}t_m)^{\mathsf T}$. Using $B_o,B_h,B_f$ for the
OO, OV and VV blocks, define for occupied $i$ in the pair
\begin{align}
 O_i &= (B_o)_{i:}+B_h n_{i:}^{\mathsf T},\\
 V_i &= (B_h)_{i:}^{\mathsf T}-n^{\mathsf T}O_i+B_f n_{i:}^{\mathsf T},
 & D &= B_f-n^{\mathsf T}B_h .
\end{align}
Here $O_i$ denotes a domain column. Differentiate each product, e.g.,
$D^v=B_f^v-(n^v)^{\mathsf T}B_h-n^{\mathsf T}B_h^v$.
Jiang Eqs.~75 and 76 then contribute to the unique doubles block
\begin{equation}
 b_{2,ij}^{75,76}=\sum_Q\big[
 V_i^v V_j^{\mathsf T}+ V_i(V_j^v)^{\mathsf T}
 + D^v t_{ij}D^{\mathsf T}+D t_{ij}(D^v)^{\mathsf T}\big]_Q .
\end{equation}
\subsection{Complete fixed-amplitude CCSD nuclear source}
Let $D_v$ differentiate the molecular inputs, lower-order solutions and
interpair transfers on the selected branch while keeping the components of
$t_i$ and $t_{ij}$ fixed in their respective moving frames. For a multilinear
contraction $C$ (including transposes and permutations), the implemented rule is
\begin{equation}
 D_v C(X_1,\ldots,X_n)=\sum_{k=1}^n
 C(X_1,\ldots,D_v X_k,\ldots,X_n).
\end{equation}
For $p=ij$, put $s_{pq}=Y_p^{\mathsf T} S Y_q$, $n_m^p=s_{p,mm}t_m$,
$u_p=2t_p-t_p^{\mathsf T}$, $L_p=2K_p-K_p^{\mathsf T}$ and
$P_{pq}(X)=s_{pq}Xs_{qp}$.
Then $D_v n_m^p=(D_v s_{p,mm})t_m$.
The additional term $s_{p,mm}t_m^v$ belongs to the amplitude Jacobian action.

For each target auxiliary row $Q$, let $h^Q$ be its occupied--virtual factor,
and $O_i^Q$ the dressed occupied column defined above. Jiang Eq.~82 gives
\begin{align}
 \beta_{p}&=\sum_Q[O_i^Q(O_j^Q)^{\mathsf T}+h^Qt_p(h^Q)^{\mathsf T}],\\
 b_{2,p}^{B}&=\sum_{kl}\big[(D_v\beta_{p,kl})P_{p,kl}(t_{kl})
       +\beta_{p,kl}D_v P_{p,kl}(t_{kl})\big].
\end{align}
For Eqs.~83--84 the occupied sums below follow the native target domain.
Write $J_{lc}=(ij|lc)$, $K_{lc}=(il|jc)$ and $V^i_{bac}=(ib|ac)$ in
the target fit. In matrix notation $n$ stacks the rows $(n_l^p)^{\mathsf T}$:
\begin{align}
 \gamma_p={}&-n^{\mathsf T}J+\sum_b(n_j^p)_b V^i_{b::}
 -\sum_l n_l^p(s_{p,il}K_{il}^{\mathsf T}n_j^{il})^{\mathsf T}
 -\tfrac12\sum_l s_{p,lj}t_{lj}s_{lj,il}K_{il}s_{il,p},\\
 (\delta_p)_{ac}={}&-[n^{\mathsf T}(2K-J)]_{ac}
 +\sum_b(n_i^p)_b(2V^j_{cab}-V^j_{bca})\notag\\
 &-\sum_l (n_l^p)_a[(n_i^{lj})^{\mathsf T}L_{lj}s_{lj,p}]_c
 +\tfrac12\sum_l[s_{p,il}u_{il}s_{il,lj}L_{lj}s_{lj,p}]_{ac}.
\end{align}
The tangent uses $D_v$ on every nonconstant operand. Use $J^{ijk},K^{ijk}$
for the native mixed integrals with their second virtual axis in $kj$
($jk$ shares that frame). The ordered terms and their permutation are
\begin{align}
 C_{ij}&=-\sum_k(J^{ijk}+s_{p,ki}\gamma_{ki}s_{ki,kj})t_{kj}^{\mathsf T}s_{kj,p},\\
 D_{ij}&=\tfrac12\sum_k(2K^{ijk}-J^{ijk}+s_{p,ik}\delta_{ik}s_{ik,jk})
                 u_{jk}^{\mathsf T}s_{jk,p},\\
 M_{ij}&=\tfrac12 C_{ij}+C_{ij}^{\mathsf T}+D_{ij},&
 b_{2,p}^{C,D}&=D_v(M_{ij}+M_{ji}^{\mathsf T}).
\end{align}
With $F''$ the fully dressed and doubles-corrected Fock intermediate of the
right equations (in its indicated frame), the last doubles terms are
\begin{align}
 E_p &=t_p(F''_p)^{\mathsf T}+F''_pt_p,&
 G_{ij}&=-\sum_k F''_{kj}P_{p,ik}(t_{ik}),\\
 b_{2,p}^{E,G}&=D_v(E_p+G_{ij}+G_{ji}^{\mathsf T}).
\end{align}
For singles let $A_i,B_i,C_i$ be the native ordered contractions in Jiang
Eqs.~88--90 expressed in the diagonal frame $ii$. In the notation of Eq.~87,
\begin{equation}
 b_{1,i}=D_v(F^{\rm dressed}_{ai}+A_i+B_i+C_i),\qquad
 b_{2,p}=b_{2,p}^{75,76}+b_{2,p}^{B}+b_{2,p}^{C,D}+b_{2,p}^{E,G}.
\end{equation}
Here $F^{\rm dressed}_{ai}$ denotes the vector over $a$. Its derivative
includes the moving occupied and virtual legs and all projected singles in
Eqs.~96 and 98--101. Each bare Fock term occurs once.
For the coupled amplitude vector $t=(t_1,t_2)$, $R_t t^v=-b^v$.
Equations~\eqref{eq:curvature} and \eqref{eq:column} explain why stationary
curvature requires first-order electronic relaxation but no second-order
amplitude derivatives; in the symmetric expression the reference multipliers
suffice.
\subsection{Coupled right amplitude response in the transformed representation}
For the Jacobian action all orbital frames, fitted factors, overlap maps and
rank choices stay at their reference values. Define
$P_x(s)=R(t_1+s x_1,t_2+s x_2)$ for the ordinary CCSD residual.
The projected singles and the occupied dressed factor are linear in the
amplitudes; the dressed virtual leg of the direct integral can be quadratic.
Its product with the other virtual leg is at most quartic. The virtual
ladder is at most cubic. The $\beta,\gamma,\delta$ terms are at most quadratic
before multiplication by a doubles amplitude. The dressed Fock residual
terms also contain at most four amplitude factors. Pair transfers are linear
with constant coefficients in this action. Thus
$P_x(s)=\sum_{m=0}^4 c_m s^m$, with $c_1=Jx$. For any nonzero $a$,
\begin{align}
 O_x(a)&=\tfrac12[P_x(a)-P_x(-a)]=a c_1+a^3 c_3,\\
 Jx&=\frac{8O_x(a)-O_x(2a)}{6a}.
\end{align}
The implementation sets $a\|x\|_\infty=1/8$ and returns zero directly for
$x=0$. Extraction is exact for this polynomial in exact arithmetic.
It uses four native residual evaluations per nonzero action. All evaluations
are at one nuclear geometry.

For the real RHF pair symmetry $x_{ji}=x_{ij}^{\mathsf T}$, the metric is
\begin{equation}
 \langle z,x\rangle_W=\sum_i z_i^{\mathsf T}x_i
  +\sum_{i<j}z_{ij}:x_{ij}+\tfrac12\sum_i z_{ii}:x_{ii},
 \qquad \langle z,Jx\rangle_W=\langle J^{\dagger_W}z,x\rangle_W .
\end{equation}
The native Lambda reverse sweep supplies $J^{\dagger_W}z$ by omitting the
additive CC energy and triples source. The doubles of a diagonal occupied
pair are symmetric. Store its upper triangle with unit weight off the
diagonal and components divided by $\sqrt{2}$ on the diagonal, so the
ordinary packed dot product implements $W$.
Preconditioned GMRES solves $Jt^v=-b^v$. Its diagonal preconditioner uses
$F_{aa}^{ii}-F_{ii}$ for singles and
$F_{aa}^{ij}+F_{bb}^{ij}-F_{ii}-F_{jj}$ for doubles in the reference
semicanonical pair frames. The unpreconditioned residual is checked at exit.
\end{document}
